%Główny plik Pracy dyplomowej

\documentclass[
    left=2.5cm,         
    right=2.5cm,        
    top=2.5cm,          
    bottom=3cm,         
    bindingoffset=6mm,  
    nohyphenation=false %Sprawdźić czy jest do wyłączenia czy nie!!!!!!!!!!!!!!
]{technical/eiti-thesis}

\langpol                              % Język Polski
\graphicspath{{figures/}}             % ścieżka katalogu z grafikami.
\addbibresource{bibliografia.bib}     % Plik z bibliografią

\begin{document}

%--------------------------------------
% Strona tytułowa
%--------------------------------------
\EngineerThesis
\instytut{Instytut Systemów Elektronicznych}
\kierunek{Elektronika}
\specjalnosc{Elektronika i fotonika}
\title{
    Opracowanie generatora impulsów do pomiarów parametrów linii transmisyjnych
}
\engtitle{
    Development of a Pulse Generator for Transmission Line Parameter Measurements
}
\author{Piotr Budny}
\album{324985}
\promotor{Dr inż. Maciej Urbański}
\date{2026}
\maketitle

%--------------------------------------
% Streszczenie po polsku
%--------------------------------------
\cleardoublepage 
\streszczenie \input{technical/abstractPOL}
\slowakluczowe XXX, XXX, XXX

%--------------------------------------
% Streszczenie po angielsku
%--------------------------------------
\newpage
\abstract \input{technical/abstractENG}
\keywords XXX, XXX, XXX

%--------------------------------------
% Oświadczenie o autorstwie
%--------------------------------------
\cleardoublepage  
\pagestyle{plain}
\makeauthorship

%--------------------------------------
% Spis treści
%--------------------------------------
\cleardoublepage 
\tableofcontents

%--------------------------------------
% Rozdziały
%--------------------------------------
\cleardoublepage 
\pagestyle{headings}

\input{1-wstep}         
\newpage
\section{Podstawy teoretyczne}

\subsection{Podstawy fizyczne}
\subsubsection{Impedancja}
\subsubsection{Fala}
\subsubsection{Zjawisko odbicia fali}
\subsubsection{Znajdowanie miejsca odbicia fali}
\subsubsection{Określanie charakterystyki odbicia}

\subsection{Podstawy generacji impulsów}
\subsubsection{Czas narastania i częstotliwość}
\subsubsection{Charakterystyka impulsów}
\subsubsection{Sposoby generowania impulsów}    
\newpage
\section{Rozwiązania dostępne na rynku}

\subsection{Rozwiązania profesjonalne}
\subsubsection{Tytuł rozwiązania nr 1}
\subsubsection{Tytuł rozwiązania nr 2}
\subsubsection{Tytuł rozwiązania nr 3}

\subsection{Rozwiązania Amatorskie}
\subsubsection{Tytuł rozwiązania nr 1}
\subsubsection{Tytuł rozwiązania nr 2}
\subsubsection{Tytuł rozwiązania nr 3}    
\newpage
\section{*Wybrane rozwiązanie*} %Jakoś inaczej może?

\subsection{*Opis działania wybranego rozwiązania*}
\subsection{Zalety i Wady rozwiązania}
\subsection{Wymagania projektowe}
    
\newpage
\section{*Projekt reflektometru*} %Architektura reflektometru/Budowa reflektometru/??

\subsection{}    
\input{6-Oprogramowanie}    
\newpage
\section{Testy}

\subsection{Testy uruchomieniowe}
\subsection{Testy funkcjonalne}
\subsection{Testy oprogramowania}
\subsection{Testy warunków granicznych}
    
\input{8-Podsumowanie}    

%--------------------------------------------
% Literatura
%--------------------------------------------
\cleardoublepage 
\printbibliography

%--------------------------------------------
% Spisy (opcjonalne)
%--------------------------------------------
\newpage
\pagestyle{plain}

% Wykaz symboli i skrótów.
\vspace{0.8cm}
\acronymlist
%\acronym{EiTI}{Wydział Elektroniki i Technik Informacyjnych}


\listoffigurestoc     % Spis rysunków.
\vspace{1cm}          
\listoftablestoc      % Spis tabel.
\vspace{1cm}          
\listofappendicestoc  % Spis załączników

% Załączniki
\newpage
%\appendix{Nazwa załącznika 1}





\end{document}
